\documentclass[a4paper,10pt,twocolumn]{article}

% =================================================
% ARTIKKELSTIL
% =================================================
\usepackage{ffvstil}
% Kommenter inn for å fjerne lenkefarger (ren typografi, ingen farger):
% \hypersetup{hidelinks}

% =================================================
% FYSIKK OG MATEMATIKK
% =================================================
\usepackage{amsmath}       % ikke amssymb: unicode-math leverer
                           % symbolene, og lastes begge brekker
                           % bygget (\eth already defined)
                           % mathtools lastes IKKE: se README for
                           % alternativene som er native i unicode-math
                           % ikke bm: inkompatibel med unicode-math.
                           % Fet matte er \symbf / \symbfit / \symbfup.

% =================================================
% EGNE PAKKER OG KOMMANDOER
%
% Legg inn ekstra pakker og egendefinerte kommandoer
% her. Eksempler som ofte er nyttige er kommentert ut
% nedenfor — fjern %-tegnet for å aktivere dem.
% =================================================

% --- Ekstra matematiske pakker ---
% \usepackage{slashed}        % Feynman-slash for Dirac-matriser
% \usepackage{tensor}         % Pene tensorindekser
% \usepackage{physics2}       % så \usephysicsmodule{braket,ab.legacy}
%                             % for \ket, \bra, \abs, \norm
% \usepackage{cancel}         % Overstrøket matematikk

% --- Ekstra tekst- og layoutpakker ---
% \usepackage{enumitem}       % Finkontroll over lister
% \usepackage{siunitx}        % SI-enheter: \qty{9.81}{\metre\per\second\squared}
% \usepackage{epigraph}       % Sitater ved seksjonsstart

% --- Egne kommandoer ---
% \newcommand{\R}{\mathbb{R}}
% \newcommand{\C}{\mathbb{C}}
% \newcommand{\mean}[1]{\langle #1 \rangle}

% --- Egne operatorer ---
% \DeclareMathOperator{\Tr}{Tr}
% \DeclareMathOperator{\sgn}{sgn}

% =================================================
% METADATA — fyll inn her
% =================================================
\title{ARTIKKELTITTEL\\[0.4em]
{\large Undertittel om ønskelig}}
\author{Fornavn Etternavn}
\date{\today}

% =================================================
% DOKUMENT
% =================================================
\begin{document}

\twocolumn[
  \maketitle
  \begin{ingress}
    Skriv ingressen her. Ingressen er et kort, selvstendig avsnitt i
    fet skrift som gir leseren et raskt overblikk over hva artikkelen
    handler om, hva som er det sentrale argumentet, og hvorfor det
    er interessant. To til fire setninger er vanligvis passe.
    Siteringer er tillatt også her~\cite{eksempel2024}.
  \end{ingress}
]

% =================================================
\section{Innledning}

Skriv innledningen her. Introduser problemstillingen, gi nødvendig
bakgrunn, og avslutt gjerne med en kort oversikt over artikkelens
struktur. Avsnitt markeres med innrykk --- ikke med ekstra
linjeskift.

% =================================================
\section{Seksjon to}

Matematikk settes inline som $E = mc^2$ eller som nummererte
ligninger:
\begin{equation}
  \hat{H}\psi = E\psi.
  \label{eq:eksempel}
\end{equation}
Referer til ligninger med \eqref{eq:eksempel} og til
referanser med \cite{eksempel2024}.

% -------------------------------------------------
% FIGUR — enkolonne (flyter inn i én kolonne)
% -------------------------------------------------
\begin{figure}[t]
  \centering
  % Bytt ut med din egen figur:
  % \includegraphics[width=\linewidth]{filnavn}
  \rule{0.9\linewidth}{3cm}     % plassholder
  \caption{Dette er en figurtekst for en enkolonne-figur.
           Teksten brytes automatisk og innrykkes pent under
           etiketten.}
  \label{fig:enkel}
\end{figure}

Henvis til figurer i teksten: figur~\ref{fig:enkel} viser et
eksempel. Plasseringsargumentet \texttt{[t]} ber LaTeX om å
plassere figuren øverst på en kolonne --- de vanligste
alternativene er \texttt{t} (topp), \texttt{b} (bunn),
\texttt{h} (her om mulig), og \texttt{p} (egen side).

\subsection{Underseksjon}

Bruk underseksjoner for å strukturere lengre seksjoner.

% -------------------------------------------------
% FIGUR — tokolonne (spenner over begge kolonner)
% -------------------------------------------------
\begin{figure*}[t]
  \centering
  % Bytt ut med din egen figur:
  % \includegraphics[width=0.9\linewidth]{filnavn}
  \rule{0.9\linewidth}{3.5cm}   % plassholder
  \caption{Dette er en tokolonne-figur som spenner over hele
           sidebredden. Brukes for brede illustrasjoner, plott
           med mange paneler, eller figurer som krever mye plass
           for å være lesbare. \texttt{figure*}-miljøet flyter
           kun til topp eller bunn av sider.}
  \label{fig:bred}
\end{figure*}

% -------------------------------------------------
% TABELL
% -------------------------------------------------
\begin{table}[t]
  \centering
  \caption{Et eksempel på en smakfull tabell med \texttt{booktabs}.
           Vertikale streker unngås; horisontale streker brukes
           sparsomt (topp, mellom hode og data, bunn).}
  \label{tab:eksempel}
  \begin{tabular}{lcr}
    \toprule
    Størrelse & Symbol & Verdi \\
    \midrule
    Lysets hastighet   & $c$      & $2{,}998 \times 10^8$ m/s \\
    Plancks konstant   & $\hbar$  & $1{,}055 \times 10^{-34}$ Js \\
    Elementærladning   & $e$      & $1{,}602 \times 10^{-19}$ C \\
    \bottomrule
  \end{tabular}
\end{table}

% =================================================
\section{Avslutning}

Oppsummer hovedresultatene og pek gjerne mot åpne spørsmål eller
videre arbeid.

% =================================================
% REFERANSER
% Kompileringsrekkefølge:
%   lualatex main
%   bibtex   main
%   lualatex main
%   lualatex main
% =================================================
\bibliographystyle{unsrtnat}
\bibliography{referanser}

\end{document}
